% Automatisch erzeugt aus anhangsverzeichnis_ano.csv
% Anzahl Einträge: 127

% Benötigte Pakete:
% \usepackage{enumitem}
% \usepackage{xurl}
% \usepackage{microtype} % optional

\begin{enumerate}[label=\arabic*, leftmargin=*, itemsep=0.5em, topsep=0.3em]
\item \textbf{code} \textnormal{(\textit{Verzeichnis})}\\
Codes zur Vorverarbeitung der Basisdaten und Berechnung sowie Auswertung der Indikatoren
\begin{enumerate}[label=\arabic{enumi}.\arabic*, leftmargin=*, itemsep=0.5em, topsep=0.3em]
\item \textbf{helper} \textnormal{(\textit{Verzeichnis})}\\
Hilfsfunktionen
\begin{enumerate}[label=\arabic{enumi}.\arabic{enumii}.\arabic*, leftmargin=*, itemsep=0.5em, topsep=0.3em]
\item \textbf{all\_clear.R} \textnormal{(\textit{R})}\\
Stoppen von r5r und Garbage Collection\\
\texttt{\small code/helper/all\_clear.R}
\item \textbf{check\_gtfs\_discont.R} \textnormal{(\textit{R})}\\
Prüfen eines GTFS-Feeds auf Diskontinuitäten\\
\texttt{\small code/helper/check\_gtfs\_discont.R}
\item \textbf{mode\_value.R} \textnormal{(\textit{R})}\\
Berechnen des Modus einer Variable\\
\texttt{\small code/helper/mode\_value.R}
\item \textbf{pointgrid\_fun.R} \textnormal{(\textit{R})}\\
Erstellung einer r5r-kompatiblen Tabelle aus einem Geogitter\\
\texttt{\small code/helper/pointgrid\_fun.R}
\item \textbf{pois\_fun.R} \textnormal{(\textit{R})}\\
Erstellung von r5r-kompatiblen Tabellen für Start- und Zielpunkte\\
\texttt{\small code/helper/pois\_fun.R}
\item \textbf{polygrid\_fun.R} \textnormal{(\textit{R})}\\
Erstellung eines Geogitters\\
\texttt{\small code/helper/polygrid\_fun.R}
\item \textbf{school\_holidays.R} \textnormal{(\textit{R})}\\
Abfrage von Schulferien per API\\
\texttt{\small code/helper/school\_holidays.R}
\item \textbf{nonholiday\_normdays\_cutoff.txt} \textnormal{(\textit{Text})}\\
Normtage im DELFI-Feed (Output von 03\_find\_valid\_dates)\\
\texttt{\small code/temp/nonholiday\_normdays\_cutoff.txt}
\item \textbf{nonholiday\_weekdays\_cutoff.txt} \textnormal{(\textit{Text})}\\
Werktage im DELFI-Feed (Output von 03\_find\_valid\_dates)\\
\texttt{\small code/temp/nonholiday\_weekdays\_cutoff.txt}
\end{enumerate}
\item \textbf{01\_osm\_extract.R} \textnormal{(\textit{R})}\\
Erzeugen von Grundlagenlayern (Gemeinden, Kreise)
OSM-Zuschnitt auf Untersuchungsgebiet
Vorverarbeitung Zentrales Haltestellenverzeichnis\\
\texttt{\small code/01\_osm\_extract.R}
\item \textbf{02\_de\_gtfs\_cleaning.R} \textnormal{(\textit{R})}\\
Zuschnitt und Vorverarbeitung des DELFI-GTFS\\
\texttt{\small code/02\_de\_gtfs\_cleaning.R}
\item \textbf{03\_find\_valid\_dates.R} \textnormal{(\textit{R})}\\
Identifizieren von validen Analysetagen anhand von Ferien und Feiertagen sowie eventuell fehlenden Daten\\
\texttt{\small code/03\_find\_valid\_dates.R}
\item \textbf{04\_stop\_frequencies\_de\_gtfs.R} \textnormal{(\textit{R})}\\
Berechnung der stündlichen Bedienungsqualitätswerte\\
\texttt{\small code/04\_stop\_frequencies\_de\_gtfs.R}
\item \textbf{05\_bq\_summaries.R} \textnormal{(\textit{R})}\\
Zusammenfassen der BQ nach verschiedenen Zeitfenstern, Erstellung von Plots\\
\texttt{\small code/05\_bq\_summaries.R}
\item \textbf{06\_grid\_creation.R} \textnormal{(\textit{R})}\\
Erstellung eines EPSG:3035-Geogitters mit 100 m Kantenlänge und INSPIRE-konformen Gitter-Ids für das Untersuchungsgebiet\\
\texttt{\small code/06\_grid\_creation.R}
\item \textbf{07\_POI\_creation.R} \textnormal{(\textit{R})}\\
Extraktion und Kategorisierung von POI-Sätzen aus OSM-Daten für Kerndichtenschätzung in 08\_cpt\_KDE.R\\
\texttt{\small code/07\_POI\_creation.R}
\item \textbf{08\_cpt\_KDE.R} \textnormal{(\textit{R})}\\
Kerndichtenschätzung von POI je Gemeinde zur Ermittlung zentralörtlicher Cluster
Bewertung der Cluster und Festlegung zentraler Orte\\
\texttt{\small code/08\_cpt\_KDE.R}
\item \textbf{09\_grid\_zensusvalues.R} \textnormal{(\textit{R})}\\
Verbinden des in 06\_grid\_creation erzeugten Gitters mit den Bevölkerungszahlen\\
\texttt{\small code/09\_grid\_zensusvalues.R}
\item \textbf{10\_i2\_traveltimes.R} \textnormal{(\textit{R})}\\
Berechnung von Gehzeiten (Zensuszellen zu Haltestellen)\\
\texttt{\small code/10\_i2\_traveltimes.R}
\item \textbf{11\_i2\_analysis.R} \textnormal{(\textit{R})}\\
Ermitteln der Erschließungsqualitäten aus Gehzeit und erreichter Bedienungsqualität\\
\texttt{\small code/11\_i2\_analysis.R}
\item \textbf{12\_i2\_summaries.R} \textnormal{(\textit{R})}\\
Aggregieren der EQ-Werte nach Stunde, Tag, Wochentag und Gesamt\\
\texttt{\small code/12\_i2\_summaries.R}
\item \textbf{13\_i2\_togrid.R} \textnormal{(\textit{R})}\\
Erzeugen von Kartenebenen für kartographische Darstellung der EQ, Plot-Darstellungen\\
\texttt{\small code/13\_i2\_togrid.R}
\item \textbf{14\_i3\_traveltimes.R} \textnormal{(\textit{R})}\\
Berechnung von ÖPNV-Wegzeiten (Zentrale Orte zu Zensuszellen)\\
\texttt{\small code/14\_i3\_traveltimes.R}
\item \textbf{15\_i3\_shortest\_times.R} \textnormal{(\textit{R})}\\
Identifizieren kürzester Wegzeiten, Berechnung des Anbindungsindikators je Stunde\\
\texttt{\small code/15\_i3\_shortest\_times.R}
\item \textbf{16\_i3\_summaries.R} \textnormal{(\textit{R})}\\
Aggregieren des Anbindungsindikators nach Tag und Stunde, erzeugen von Kartenebenen und Plots\\
\texttt{\small code/16\_i3\_summaries.R}
\item \textbf{dataenv.R} \textnormal{(\textit{R})}\\
Einträge des Fahrplandatums und ZHV-Datums zur Abfrage in diversen Skripten, Farbpaletten\\
\texttt{\small code/dataenv.R}
\item \textbf{erschließung\_mat\_long\_numeric\_stingy.csv} \textnormal{(\textit{CSV})}\\
Zuordnung von BQ/Gehzeit-Kombinationen zu Erschließungsqualitäten\\
\texttt{\small code/erschließung\_mat\_long\_numeric\_stingy.csv}
\item \textbf{figures\_stop\_frequencies.R} \textnormal{(\textit{R})}\\
Frühere Version von 04\_stop\_frequencies\_de\_gtfs, genutzt um Übersichtsgrafiken zur Bedienungsqualität zu erstellen.\\
\texttt{\small code/figures\_stop\_frequencies.R}
\item \textbf{impedance\_functions.R} \textnormal{(\textit{R})}\\
Erstellung der Grafik zu Impedanzfunktionen in Abschnitt 2.2.1\\
\texttt{\small code/impedance\_functions.R}
\item \textbf{kde\_parallel.R} \textnormal{(\textit{R})}\\
Alternative zu einem Teil von 08\_cpt\_KDE. Parallele Berechnung der KDE-Raster.\\
\texttt{\small code/kde\_parallel.R}
\item \textbf{kreis\_order.txt} \textnormal{(\textit{Text})}\\
Reihenfolge der Kreise nach gewichtetem Mittel der Erschließungsqualität. Output von 13\_i2\_togrid.\\
\texttt{\small code/kreis\_order.txt}
\item \textbf{mean\_frequencies.R} \textnormal{(\textit{R})}\\
Berechnung von Bedienungsqualitäten auf Grundlage täglicher, wochentäglicher und gesamter Mittel der Abfahrten pro Stunde.\\
\texttt{\small code/mean\_frequencies.R}
\item \textbf{osmconf\_cpt.txt} \textnormal{(\textit{Text})}\\
Konfiguration von osmextract zum Extrahieren von POI aus OSM-Daten.\\
\texttt{\small code/osmconf\_cpt.txt}
\item \textbf{stop\_frequencies\_methods.R} \textnormal{(\textit{R})}\\
Frühere Version von 04\_stop\_frequencies\_de\_gtfs, genutzt um Grafiken zur stichprobenartigen Überprüfung der Ergebnisse zu erstellen.
Im weiteren Verlauf nicht aktualisiert, daher nicht funktionsfähig.\\
\texttt{\small code/stop\_frequencies\_methods.R}
\item \textbf{tags\_filter.cmd} \textnormal{(\textit{CMD})}\\
osmium-Ausdruck zum Vorfiltern des OSM-Datensatzes\\
\texttt{\small code/tags\_filter.cmd}
\end{enumerate}
\item \textbf{document} \textnormal{(\textit{Verzeichnis})}\\
LaTeX-Arbeitsverzeichnis
\begin{enumerate}[label=\arabic{enumi}.\arabic*, leftmargin=*, itemsep=0.5em, topsep=0.3em]
\item \textbf{figures} \textnormal{(\textit{Verzeichnis})}\\
Abbildungen im Dokument
\item \textbf{tables} \textnormal{(\textit{Verzeichnis})}\\
Tabellen im Dokument
\item \textbf{routingindicators.pdf} \textnormal{(\textit{PDF})}\\
Dieses Dokument\\
\texttt{\small document/routingindicators.pdf}
\item \textbf{routingindicators.tex} \textnormal{(\textit{LaTeX})}\\
TeX-Datei zu diesem Dokument\\
\texttt{\small document/routingindicators.tex}
\end{enumerate}
\item \textbf{geodata} \textnormal{(\textit{Verzeichnis})}\\
Geodaten, hauptsächlich vorverarbeitete Input-Dateien
\begin{enumerate}[label=\arabic{enumi}.\arabic*, leftmargin=*, itemsep=0.5em, topsep=0.3em]
\item \textbf{zHV\_aktuell\_csv.2026-05-21.zip} \textnormal{(\textit{ZIP})}\\
Auszug des zentralen Haltestellenverzeichnis vom 21.05.2026\\
\texttt{\small geodata/zHV\_aktuell\_csv.2026-05-21.zip}
\item \textbf{dvg1nw.gpkg} \textnormal{(\textit{GeoPackage})}\\
Verwaltungsgrenzen und Untersuchungsgebiet\\
\texttt{\small geodata/dvg1nw.gpkg}
\begin{enumerate}[label=\arabic{enumi}.\arabic{enumii}.\arabic*, leftmargin=*, itemsep=0.5em, topsep=0.3em]
\item \textbf{regbez10kmbuffer} \textnormal{(\textit{Layer})}\\
Untersuchungsgebiet mit 10 km Puffer
\item \textbf{regbez25kmbuffer} \textnormal{(\textit{Layer})}\\
Untersuchungsgebiet mit 25 km Puffer
\item \textbf{gemeinden\_regbez\_25km} \textnormal{(\textit{Layer})}\\
Gemeindengrenzen des Untersuchungsgebietes mit 25 km Puffer (VG250)
\item \textbf{gemeinden\_regbez\_vg250} \textnormal{(\textit{Layer})}\\
Gemeindengrenzen des Untersuchungsgebietes (VG250)
\item \textbf{kreise\_regbez\_vg250} \textnormal{(\textit{Layer})}\\
Kreisgrenzen des Untersuchungsgebietes (VG250)
\end{enumerate}
\item \textbf{gtfs\_shapes.gpkg} \textnormal{(\textit{GeoPackage})}\\
Mit pfaedle georeferenzierte shapes des DELFI-GTFS vom 18.05.2026\\
\texttt{\small geodata/gtfs\_shapes.gpkg}
\item \textbf{zentrale\_orte.gpkg} \textnormal{(\textit{GeoPackage})}\\
Routingziele für die Indikatorberechnungen\\
\texttt{\small geodata/zentrale\_orte.gpkg}
\begin{enumerate}[label=\arabic{enumi}.\arabic{enumii}.\arabic*, leftmargin=*, itemsep=0.5em, topsep=0.3em]
\item \textbf{zhv\_20260521} \textnormal{(\textit{Layer})}\\
Zuschnitt des Zentralen Haltestellenverzeichnis
\item \textbf{zentrale\_orte\_rlp} \textnormal{(\textit{Layer})}\\
Zentrale Orte in Rheinland-Pfalz, auf Grundlage der Regionalpläne
\item \textbf{zentrale\_orte} \textnormal{(\textit{Layer})}\\
Initiale Ergebnisse der Kerndichtenschätzung
\item \textbf{zentrale\_orte\_oz\_manual} \textnormal{(\textit{Layer})}\\
Überarbeitete Zentrale Orte auf Grundlage der Zentrenkonzepte der OZ
\item \textbf{zentraleOrte\_routingdestinations\_zentrenkonzept} \textnormal{(\textit{Layer})}\\
Zusammengefasste Zentrale Orte (Zielorte Anbindungsindikator)
\end{enumerate}
\item \textbf{pois.gpkg} \textnormal{(\textit{GeoPackage})}\\
POI für die Kerndichtenschätzung\\
\texttt{\small geodata/pois.gpkg}
\begin{enumerate}[label=\arabic{enumi}.\arabic{enumii}.\arabic*, leftmargin=*, itemsep=0.5em, topsep=0.3em]
\item \textbf{zo\_POI\_flex\_260521} \textnormal{(\textit{Layer})}\\
POI Basisfunktionen
\item \textbf{zo\_POI\_flex\_opt260521} \textnormal{(\textit{Layer})}\\
POI Erweiterte Funktionen
\item \textbf{poi\_flex\_tt\_15} \textnormal{(\textit{Layer})}\\
POI Basisfunktionen mit Reisezeiten von Zentralen Orten zur Bewertung
\item \textbf{poi\_flex\_opt\_tt\_15} \textnormal{(\textit{Layer})}\\
POI Erweiterte Funktionen mit Reisezeiten von Zentralen Orten zur Bewertung
\end{enumerate}
\item \textbf{zensus.gpkg} \textnormal{(\textit{GeoPackage})}\\
Zensusgitter (bewohnte Zellen)\\
\texttt{\small geodata/zensus.gpkg}
\item \textbf{zentrale\_orte\_area.gpkg} \textnormal{(\textit{GeoPackage})}\\
Vektorisierte und klassifizierte Ergebnisse der Kerndichtenschätzung\\
\texttt{\small geodata/zentrale\_orte\_area.gpkg}
\begin{enumerate}[label=\arabic{enumi}.\arabic{enumii}.\arabic*, leftmargin=*, itemsep=0.5em, topsep=0.3em]
\item \textbf{areas\_flex\_catcount} \textnormal{(\textit{Layer})}\\
Dichteflächen, ermittelt aus den Ergebnissen der KDE, mit Anzahl der vorhandenen POI nach Kategorie
\item \textbf{acc\_scoring\_centroids} \textnormal{(\textit{Layer})}\\
Nach der in Abschnitt 3.3.2.1 beschriebenen Methode bewertete Zentroide der Dichteflächen
\end{enumerate}
\end{enumerate}
\item \textbf{GIS} \textnormal{(\textit{Verzeichnis})}\\
QGIS-Projekt
\begin{enumerate}[label=\arabic{enumi}.\arabic*, leftmargin=*, itemsep=0.5em, topsep=0.3em]
\item \textbf{styles} \textnormal{(\textit{Verzeichnis})}\\
Projektstile
\end{enumerate}
\item \textbf{output} \textnormal{(\textit{Verzeichnis})}\\
Outputs
\begin{enumerate}[label=\arabic{enumi}.\arabic*, leftmargin=*, itemsep=0.5em, topsep=0.3em]
\item \textbf{daily\_eq} \textnormal{(\textit{Verzeichnis})}
\begin{enumerate}[label=\arabic{enumi}.\arabic{enumii}.\arabic*, leftmargin=*, itemsep=0.5em, topsep=0.3em]
\item \textbf{results\_full\_day.parquet} \textnormal{(\textit{PARQUET})}\\
Tagesmittel der Erschließungsqualität\\
\texttt{\small output/daily\_eq/results\_full\_day.parquet}
\end{enumerate}
\item \textbf{hourly\_eq} \textnormal{(\textit{Verzeichnis})}
\begin{enumerate}[label=\arabic{enumi}.\arabic{enumii}.\arabic*, leftmargin=*, itemsep=0.5em, topsep=0.3em]
\item \textbf{results\_full\_hour.parquet} \textnormal{(\textit{PARQUET})}\\
Stündliche Erschließungsqualität\\
\texttt{\small output/hourly\_eq/results\_full\_hour.parquet}
\end{enumerate}
\item \textbf{hourmean\_eq} \textnormal{(\textit{Verzeichnis})}
\begin{enumerate}[label=\arabic{enumi}.\arabic{enumii}.\arabic*, leftmargin=*, itemsep=0.5em, topsep=0.3em]
\item \textbf{results\_full\_hourmean.parquet} \textnormal{(\textit{PARQUET})}\\
Stundenmittel der Erschließungsqualität\\
\texttt{\small output/hourmean\_eq/results\_full\_hourmean.parquet}
\end{enumerate}
\item \textbf{total\_eq} \textnormal{(\textit{Verzeichnis})}
\begin{enumerate}[label=\arabic{enumi}.\arabic{enumii}.\arabic*, leftmargin=*, itemsep=0.5em, topsep=0.3em]
\item \textbf{eq\_2026-06-26.fst} \textnormal{(\textit{FST})}\\
Mittel der Erschließungsqualität über den gesamten Zeitraum\\
\texttt{\small output/total\_eq/eq\_2026-06-26.fst}
\end{enumerate}
\item \textbf{daily\_20260518\_weekday\_2026-05-04\_2026-06-26\_.fst} \textnormal{(\textit{FST})}\\
Bedienungsqualitäten der Haltestellen auf Grundlage des täglichen Mittels der Abfahrten pro Stunde (Output von mean\_frequencies)\\
\texttt{\small output/daily\_20260518\_weekday\_2026-05-04\_2026-06-26\_.fst}
\item \textbf{dow\_20260518\_weekday\_2026-05-04\_2026-06-26\_.fst} \textnormal{(\textit{FST})}\\
Bedienungsqualitäten der Haltestellen auf Grundlage des wochentäglichen Mittels der Abfahrten pro Stunde (Output von mean\_frequencies)\\
\texttt{\small output/dow\_20260518\_weekday\_2026-05-04\_2026-06-26\_.fst}
\item \textbf{hourmean\_20260518\_weekday\_2026-05-04\_2026-06-26\_.fst} \textnormal{(\textit{FST})}\\
Bedienungsqualitäten der Haltestellen auf Grundlage des stündlichen Mittels der Abfahrten pro Stunde (Output von mean\_frequencies)\\
\texttt{\small output/hourmean\_20260518\_weekday\_2026-05-04\_2026-06-26\_.fst}
\item \textbf{totalmean\_20260518\_weekday\_2026-05-04\_2026-06-26\_.fst} \textnormal{(\textit{FST})}\\
Bedienungsqualitäten der Haltestellen auf Grundlage des Mittels der Abfahrten pro Stunde über den gesamten Zeitraum (Output von mean\_frequencies)\\
\texttt{\small output/totalmean\_20260518\_weekday\_2026-05-04\_2026-06-26\_.fst}
\item \textbf{hourly\_20260518\_weekday\_2026-05-04\_2026-06-26\_.fst} \textnormal{(\textit{FST})}\\
Bedienungsqualitäten der Haltestellen auf Grundlage der stündlichen Abfahrten pro Stunde (Output von 04\_stop\_frequencies\_de\_gtfs, Grundlage für 11\_i2\_analysis und EQ-Werte)\\
\texttt{\small output/hourly\_20260518\_weekday\_2026-05-04\_2026-06-26\_.fst}
\item \textbf{i3\_shortest\_hourly} \textnormal{(\textit{Verzeichnis})}
\begin{enumerate}[label=\arabic{enumi}.\arabic{enumii}.\arabic*, leftmargin=*, itemsep=0.5em, topsep=0.3em]
\item \textbf{captured\_mintt.fst} \textnormal{(\textit{FST})}\\
Zusammengefasste kürzeste Reisezeiten zu zentralen Orten (Output von 15\_i3\_shortest\_times)\\
\texttt{\small output/i3\_shortest\_hourly/captured\_mintt.fst}
\end{enumerate}
\item \textbf{i3\_ttm\_hourly} \textnormal{(\textit{Verzeichnis})}\\
360 ÖPNV-Reisezeitmatrizen für jede untersuchte Stunde (Output von 14\_i3\_traveltimes)
Benannt als i3\_ttm\_yyyy-MM-dd\_HH.fst
\item \textbf{kderasters} \textnormal{(\textit{Verzeichnis})}\\
177 KDE-Raster (Output von 08\_cpt\_KDE)
Benannt als kde\_\{Gemeindeschlüssel\}\_flex.tif
\item \textbf{walk\_20min\_zensus\_stops.csv} \textnormal{(\textit{CSV})}\\
Fußweg-Reisezeitmatrix zu Haltestellen (Output von 10\_i2\_traveltimes)\\
\texttt{\small output/walk\_20min\_zensus\_stops.csv}
\item \textbf{Bedienungsqualität.gpkg} \textnormal{(\textit{GeoPackage})}\\
Haltestellen mit Variabilitätswerten der BQ im und um das Untersuchungsgebiet, Grundlage für 10\_i2\_traveltimes\\
\texttt{\small output/Bedienungsqualität.gpkg}
\end{enumerate}
\item \textbf{appendix} \textnormal{(\textit{Verzeichnis})}\\
Vollständige Indikatorergebnisse auf Gemeinde- und Gittereebene
\begin{enumerate}[label=\arabic{enumi}.\arabic*, leftmargin=*, itemsep=0.5em, topsep=0.3em]
\item \textbf{figures} \textnormal{(\textit{Verzeichnis})}\\
Plots
\begin{enumerate}[label=\arabic{enumi}.\arabic{enumii}.\arabic*, leftmargin=*, itemsep=0.5em, topsep=0.3em]
\item \textbf{percbq} \textnormal{(\textit{Verzeichnis})}\\
Bedienungsqualität
\begin{enumerate}[label=\arabic{enumi}.\arabic{enumii}.\arabic{enumiii}.\arabic*, leftmargin=*, itemsep=0.5em, topsep=0.3em]
\item \textbf{dately\_gem} \textnormal{(\textit{Verzeichnis})}\\
99 Vektorgrafiken zur Entwicklung des Tages-50\%-Quantils der Bedienungsqualität im Untersuchungszeitraum je Gemeinde
\item \textbf{hourly\_gem} \textnormal{(\textit{Verzeichnis})}\\
99 Vektorgrafiken zur Entwicklung des Stunden-50\%-Quantils der Bedienungsqualität im Untersuchungszeitraum je Gemeinde
\item \textbf{Gemeinde\_bq\_per\_date.svg} \textnormal{(\textit{SVG})}\\
Übersichtsgrafik zur Entwicklung des Tages-50\%-Quantils der Bedienungsqualität der Gemeinden\\
\texttt{\small appendix/figures/percbq/Gemeinde\_bq\_per\_date.svg}
\item \textbf{Gemeinde\_bq\_per\_hour.svg} \textnormal{(\textit{SVG})}\\
Übersichtsgrafik zur Entwicklung des Stunden-50\%-Quantils der Bedienungsqualität der Gemeinden\\
\texttt{\small appendix/figures/percbq/Gemeinde\_bq\_per\_hour.svg}
\item \textbf{Kreis\_bq\_per\_date.svg} \textnormal{(\textit{SVG})}\\
Übersichtsgrafik zur Entwicklung des Tages-50\%-Quantils der Bedienungsqualität der Kreise\\
\texttt{\small appendix/figures/percbq/Kreis\_bq\_per\_date.svg}
\item \textbf{Kreis\_bq\_per\_hour.svg} \textnormal{(\textit{SVG})}\\
Übersichtsgrafik zur Entwicklung des Stunden-50\%-Quantils der Bedienungsqualität der Kreise\\
\texttt{\small appendix/figures/percbq/Kreis\_bq\_per\_hour.svg}
\item \textbf{RegioStar7\_bq\_per\_date.svg} \textnormal{(\textit{SVG})}\\
Übersichtsgrafik zur Entwicklung des Tages-50\%-Quantils der Bedienungsqualität der RegioStaR7-Klassen\\
\texttt{\small appendix/figures/percbq/RegioStar7\_bq\_per\_date.svg}
\item \textbf{RegioStaR7\_bq\_per\_hour.svg} \textnormal{(\textit{SVG})}\\
Übersichtsgrafik zur Entwicklung des Stunden-50\%-Quantils der Bedienungsqualität der RegioStaR7-Klassen\\
\texttt{\small appendix/figures/percbq/RegioStaR7\_bq\_per\_hour.svg}
\end{enumerate}
\item \textbf{perceq} \textnormal{(\textit{Verzeichnis})}\\
Erschließungsqualität
\begin{enumerate}[label=\arabic{enumi}.\arabic{enumii}.\arabic{enumiii}.\arabic*, leftmargin=*, itemsep=0.5em, topsep=0.3em]
\item \textbf{hourly} \textnormal{(\textit{Verzeichnis})}\\
99 Vektorgrafiken zur Entwicklung des Stunden-50\%-Quantils der Erschließungsqualität im Untersuchungszeitraum je Gemeinde
\item \textbf{Gemeinde\_eq\_per\_hour.svg} \textnormal{(\textit{SVG})}\\
Übersichtsgrafik zur Entwicklung des Stunden-50\%-Quantils der Erschließungsqualität der Gemeinden\\
\texttt{\small appendix/figures/perceq/Gemeinde\_eq\_per\_hour.svg}
\end{enumerate}
\item \textbf{perci3} \textnormal{(\textit{Verzeichnis})}\\
Anbindungsindikator
\begin{enumerate}[label=\arabic{enumi}.\arabic{enumii}.\arabic{enumiii}.\arabic*, leftmargin=*, itemsep=0.5em, topsep=0.3em]
\item \textbf{hourly} \textnormal{(\textit{Verzeichnis})}\\
Stündliche Entwicklung
\item \textbf{gz} \textnormal{(\textit{Verzeichnis})}\\
99 Vektorgrafiken zur stündlichen Entwicklung der Erschließung zentraler Orte je Gemeinde
\item \textbf{hourly\_tt} \textnormal{(\textit{Verzeichnis})}\\
99 Vektorgrafiken zur stündlichen Entwicklung der Reisezeit in zentrale Orte je Gemeinde
\end{enumerate}
\item \textbf{stop\_routes} \textnormal{(\textit{Verzeichnis})}
\item \textbf{stopvar} \textnormal{(\textit{Verzeichnis})}
\end{enumerate}
\item \textbf{reports} \textnormal{(\textit{Verzeichnis})}\\
Detailkarten der Indikatoren
\begin{enumerate}[label=\arabic{enumi}.\arabic{enumii}.\arabic*, leftmargin=*, itemsep=0.5em, topsep=0.3em]
\item \textbf{BQ\_Atlas.pdf} \textnormal{(\textit{PDF})}\\
Darstellung der Haltestellen je Gemeinde mit Bedienungsqualität und Variabilitätsmaßen. 1 Seite je Gemeinde.\\
\texttt{\small appendix/reports/BQ\_Atlas.pdf}
\item \textbf{EQ\_Atlas.pdf} \textnormal{(\textit{PDF})}\\
Darstellung der Zensuszellen je Gemeinde mit Erschließungsqualität und Variabilitätsmaßen. 1 Seite je Gemeinde.\\
\texttt{\small appendix/reports/EQ\_Atlas.pdf}
\item \textbf{Erschließungsqualität.pdf} \textnormal{(\textit{PDF})}\\
Größerformatige Darstellung nur der Erschhlließungsqualität. 1 Seite je Gemeinde.\\
\texttt{\small appendix/reports/Erschließungsqualität.pdf}
\item \textbf{Indikator\_03\_Atlas.pdf} \textnormal{(\textit{PDF})}\\
Darstelllung der Erreichbarkeit zentraler Orte je Gemeinde mit Variabilitätsmaßen. 4 Seiten je Gemeinde.\\
\texttt{\small appendix/reports/Indikator\_03\_Atlas.pdf}
\end{enumerate}
\end{enumerate}
\item \textbf{results} \textnormal{(\textit{Verzeichnis})}\\
Aggregierte Ergebnisse
\begin{enumerate}[label=\arabic{enumi}.\arabic*, leftmargin=*, itemsep=0.5em, topsep=0.3em]
\item \textbf{i2\_summary.fst} \textnormal{(\textit{FST})}\\
Erschließungsqualität, aggregiert über gesamten Zeitraum\\
\texttt{\small results/i2\_summary.fst}
\item \textbf{i2\_summary\_by\_date.fst} \textnormal{(\textit{FST})}\\
Erschließungsqualität, aggregiert nach Datum\\
\texttt{\small results/i2\_summary\_by\_date.fst}
\item \textbf{i2\_summary\_by\_hour.fst} \textnormal{(\textit{FST})}\\
Erschließungsqualität, aggregiert nach Stunde\\
\texttt{\small results/i2\_summary\_by\_hour.fst}
\item \textbf{i2\_summary\_by\_weekday.fst} \textnormal{(\textit{FST})}\\
Erschließungsqualität, aggregiert nach Wochentag\\
\texttt{\small results/i2\_summary\_by\_weekday.fst}
\item \textbf{i3\_summary\_by\_date.fst} \textnormal{(\textit{FST})}\\
Erreichbarkeit zentraler Orte, aggregiert nach Datum\\
\texttt{\small results/i3\_summary\_by\_date.fst}
\item \textbf{i3\_summary\_by\_hour.fst} \textnormal{(\textit{FST})}\\
Erreichbarkeit zentraler Orte, aggregiert nach Stunde\\
\texttt{\small results/i3\_summary\_by\_hour.fst}
\item \textbf{indikator\_02.gpkg} \textnormal{(\textit{GeoPackage})}\\
Kartendaten Ergebnisse Erschließungsqualität, Grundlage für die PDF-Dateien in /appendix/reports\\
\texttt{\small results/indikator\_02.gpkg}
\begin{enumerate}[label=\arabic{enumi}.\arabic{enumii}.\arabic*, leftmargin=*, itemsep=0.5em, topsep=0.3em]
\item \textbf{i2aggregated\_hourly\_2026-05-04\_2026-06-26\_20260518} \textnormal{(\textit{Layer})}\\
Erschließungsqualität, aggregiert nach Gemeinde auf Grundlage der stündlichen Werte mit Bevölkerungsanteilen und Abweichung vom Soll
\item \textbf{i2grid\_hourly\_2026-05-04\_2026-06-26\_20260518} \textnormal{(\textit{Layer})}\\
Geogitter Erschließungsqualität mit Variabilitätswerten und Bevölkerungsdaten
\end{enumerate}
\item \textbf{indikator\_03.gpkg} \textnormal{(\textit{GeoPackage})}\\
Kartendaten Ergebnisse Erreichbarkeit zentraler Orte, Grundlage für die PDF-Dateien in /appendix/reports\\
\texttt{\small results/indikator\_03.gpkg}
\begin{enumerate}[label=\arabic{enumi}.\arabic{enumii}.\arabic*, leftmargin=*, itemsep=0.5em, topsep=0.3em]
\item \textbf{i3\_summary\_gem} \textnormal{(\textit{Layer})}\\
Erreichbarkeit zentraler Orte, aggregiert nach Gemeinde auf Grundlage der stündlichen Werte
\item \textbf{capture\_grid} \textnormal{(\textit{Layer})}\\
Geogitter Erreichbarkeit zentraler Orte mit Variabilitätswerten und Reisezeiten
\end{enumerate}
\item \textbf{Bedienungsqualität.gpkg} \textnormal{(\textit{GeoPackage})}\\
Zuschnitt der Ergebnisse der Bedienungsqualität auf das Untersuchungsgebiet\\
\texttt{\small results/Bedienungsqualität}
\end{enumerate}
\end{enumerate}
